\section{Positive-Diagonal Matignon Framework}\label{sec:framework}

Fix \(0<\alpha<1\) and define
\begin{equation}
 \Sigma_\alpha
 :=
 \left\{z\in\mathbb C\setminus\{0\}:|\arg z|>\frac{\alpha\pi}{2}\right\}.
\end{equation}
For real \(n\times n\) matrices,
\begin{equation}
 \mathcal F_\alpha^{(n)}
 :=
 \left\{
 A\in\mathbb R^{n\times n}:
 \spec(DA)\subset\Sigma_\alpha
 \ \text{for every positive diagonal }D
 \right\},
 \label{eq:Falpha}
\end{equation}
and let \(\mathcal D_H^{(n)}\) denote classical Hurwitz \(D\)-stability. The genuinely fractional difference class is
\begin{equation}
 \mathcal P_\alpha^{(n)}
 :=
 \mathcal F_\alpha^{(n)}\setminus\mathcal D_H^{(n)}.
 \label{eq:Palpha}
\end{equation}
Here \(D\)-stability is used only in the multiplicative matrix sense, not in the distinct control-theoretic meaning of root placement in a prescribed \(D\)-region. Matignon's criterion identifies \(\Sigma_\alpha\) as the asymptotic-stability region of a commensurate Caputo linear system \cite{Matignon1996,BrandiburGarrappaKaslik2021}.

\subsection{Projective Coordinates of the Positive-Diagonal Orbit}

Suppose \(-A\) is a strict \(P\)-matrix. Write \(p_i=-a_{ii}>0\) and \(m_I=(-1)^{|I|}\det A[I]>0\). For \(D=\diag(d_1,\ldots,d_n)\succ0\), the characteristic coefficients, projective multiplier coordinates, and normalized principal-minor invariants are
\begin{align}
 c_k(D)&=\sum_{|I|=k}m_I\prod_{i\in I}d_i,\qquad
 x_i=\frac{p_id_i}{\sum_jp_jd_j},\qquad \sum_i x_i=1,\label{eq:orbit-compact}\\
 \beta_I&=\frac{m_I}{\prod_{i\in I}p_i}\qquad(|I|\ge2).
\end{align}
Thus the positive-diagonal cone modulo common scale is exactly the open simplex. If \(a_D=\sum_ip_id_i\) and the spectral variable is scaled by \(a_D\), the orbit becomes
\begin{equation}
 z^n+z^{n-1}+B_2(x)z^{n-2}+\cdots+B_n(x),
 \qquad
 B_k(x)=\sum_{|I|=k}\beta_I\prod_{i\in I}x_i.
 \label{eq:normalized-general}
\end{equation}
This projective reduction is structural rather than a novelty claim; the principal-minor representation is classical and underlies feedback-cycle descriptions \cite{Kinkhabwala2015}.

For \(n=3\), only four nontrivial orbit invariants remain:
\begin{equation}
 \boxed{
 \beta_{12}=\frac{m_{12}}{p_1p_2},\quad
 \beta_{13}=\frac{m_{13}}{p_1p_3},\quad
 \beta_{23}=\frac{m_{23}}{p_2p_3},\quad
 \kappa=\frac{-\det A}{p_1p_2p_3}.
 }
 \label{eq:four-invariants}
\end{equation}
Section~\ref{sec:three} shows that these four quantities are sufficient to eliminate the universal positive-diagonal multiplier from the real three-dimensional Matignon problem.
